\chapter{Financial Schemas, Replication, and Connection Pooling}
\label{ch08:capstone}

For seven chapters we have worked inside the PostgreSQL engine room. We walked
8~KB slotted heap pages and buffer descriptors in Chapter~2, split Lehman--Yao
B-link trees and generalized search trees in Chapter~3, flushed write-ahead log
records and full-page images in Chapter~4, traced multiversion snapshots and
\texttt{VACUUM} in Chapter~5, reasoned about row-level locks and serializable
anomalies in Chapter~6, and read cost-based execution plans in Chapter~7. Now we
sit at the architect's desk and build a system where every one of those low-level
mechanisms has a dollar sign attached to it: a production double-entry financial
ledger served by an asynchronous Rust (\texttt{Axum} + \texttt{SQLx}) API over
PostgreSQL.

Financial systems punish vague thinking because every shortcut fails in a
specific, observable way. Store balances in binary floating point and sums of
fee postings stop matching the books to the cent. Store wall-clock timestamps
without a UTC reference and the autumn daylight-saving fallback folds two distinct
hours onto the same nominal minute. Choose random \texttt{UUIDv4} primary keys and
a sixty-gigabyte B-tree turns small inserts into an 8~KB full-page-write storm
across cold NVMe blocks. Validate invariants only in Rust handlers and a single
concurrent retry or backfill script creates money out of thin air. Leave
replication asynchronous and a primary disk failure erases transfers you already
confirmed to customers; point five thousand Tokio tasks directly at PostgreSQL and
shared-memory snapshot contention brings a thirty-two-core server to its knees.

In this final chapter, we state every architectural rule by starting with the
exact production failure it prevents, then trace how PostgreSQL's storage,
indexing, WAL, locking, and wire-protocol internals enforce the guarantee.

\begin{focusbox}
By the end of this chapter you will be able to:
\begin{itemize}
  \item Select exact storage types for money (\texttt{BIGINT} minor units versus
    base-$10{,}000$ \texttt{NUMERIC}), instants in time (\texttt{TIMESTAMPTZ}),
    and high-throughput keys (time-ordered \texttt{UUIDv7} versus \texttt{UUIDv4}).
  \item Encode domain invariants directly in the schema with \texttt{NOT NULL},
    \texttt{CHECK}, idempotency keys, and GiST temporal exclusion constraints
    (\texttt{EXCLUDE USING gist}).
  \item Implement an append-only double-entry ledger enforced at \texttt{COMMIT}
    time by a \texttt{DEFERRABLE INITIALLY DEFERRED} constraint trigger, paired
    with deadlock-free Rust (\texttt{Axum} + \texttt{SQLx}) transfer code.
  \item Partition multi-billion-row posting tables by time range for warm buffer
    locality and zero-bloat retention archival.
  \item Tune physical streaming replication across the \texttt{synchronous\_commit}
    durability ladder and stream decoded row events with logical replication.
  \item Multiplex thousands of async Rust tasks onto dozens of PostgreSQL backend
    processes through \texttt{PgBouncer} transaction pooling without breaking
    prepared statements or session state.
\end{itemize}
\end{focusbox}

\chapterline

\section{Data Types That Never Lie About Money, Time, or Identity}
\label{ch08:types}

A relational schema is a physical contract about the bit patterns allowed onto
8~KB heap pages. If a column's type rounds fractions, discards UTC context, or
scatters keys uniformly across disk, no amount of application-layer discipline can
undo the damage downstream.

\subsection{Money: Why Binary Floats Fail, and \texttt{BIGINT} vs.\ \texttt{NUMERIC}}
\label{ch08:money-types}

Suppose a payment platform stores account adjustments in PostgreSQL's 8-byte
\texttt{DOUBLE PRECISION} type (\texttt{float8}, matching \texttt{f64} in Rust and
\texttt{double} in Java or C++). A merchant collects a $\$0.10$ platform fee and a
$\$0.20$ network fee, expecting a $\$0.30$ total. Ask PostgreSQL whether the two
sides match:

\begin{lstlisting}[language=SQL]
SELECT (0.1::float8 + 0.2::float8) = 0.3::float8 AS exact_match,
       (0.1::float8 + 0.2::float8) - 0.3::float8 AS drift;
-- exact_match: false | drift: 5.551115123125783e-17
\end{lstlisting}

The comparison returns \texttt{false} because IEEE~754 binary floating point
represents numbers as $(-1)^s \times m \times 2^e$ with a 53-bit binary
significand $m$. In base~2, one tenth is an infinite repeating fraction,
$0.1_{10} = 0.0001100110011\overline{0011}_2$, just as $1/3 = 0.333\ldots_{10}$
never terminates in base~10. Rounding at bit~53 introduces an error of rough
magnitude $10^{-17}$ on a single addition. Sum ten million fee postings across
parallel \texttt{HashAggregate} workers (Chapter~7), where summation order varies
from run to run, and two identical reconciliation queries can disagree by cents.
Binary floats (\texttt{REAL} and \texttt{DOUBLE PRECISION}) are therefore strictly
forbidden for monetary quantities.

PostgreSQL gives you two exact alternatives, each suited to a clear regime:

\begin{enumerate}
  \item \textbf{\texttt{BIGINT} (\texttt{int8}) minor units.} Instead of storing
    dollars as fractions, store integer counts of the asset's smallest indivisible
    unit: cents ($10^{-2}$~USD), micro-dollars ($10^{-6}$~USD) for sub-cent API
    billing, or satoshis ($10^{-8}$~BTC). On disk, \texttt{BIGINT} occupies a
    fixed 8~bytes on the heap page, maps directly to \texttt{i64} in Rust or
    \texttt{long} in Java, and adds or compares in a single CPU integer instruction.
    Its signed range reaches $\pm(2^{63}-1) \approx \pm 9.22 \times 10^{18}$: in
    USD cents that exceeds $\$92$~quadrillion; in micro-dollars it exceeds
    $\$9.22$~trillion. Whenever you store minor units, store the asset's ISO code
    or integer exponent alongside it so nobody adds Japanese yen ($10^{0}$) to US
    cents ($10^{-2}$) or Kuwaiti dinar fils ($10^{-3}$).
  \item \textbf{\texttt{NUMERIC(p, s)} (equivalently \texttt{DECIMAL(p, s)}).}
    Foreign-exchange rates (\texttt{1.08437525}), fractional equity shares
    (\texttt{0.00418275} shares), and multi-asset bond accruals often combine
    precision $p$ up to dozens of digits with varying decimal scales $s$, or
    multiply two high-precision figures where an intermediate 64-bit integer would
    overflow. Internally (\texttt{NumericVar} in PostgreSQL's \texttt{numeric.c}),
    \texttt{NUMERIC} is a variable-length (\texttt{varlena}) structure that packs
    \emph{four base-10 digits into each 16-bit integer word} (base~$10{,}000$),
    paired with a compact header storing the base-$10{,}000$ weight, sign, and
    display scale. Because every power of $10^{-1}$ terminates in base~$10{,}000$,
    \texttt{0.1::numeric + 0.2::numeric = 0.3::numeric} is exact to every declared
    digit. The tradeoff is CPU and space: \texttt{NUMERIC} carries a variable-length
    header (typically 8--16~bytes on disk for moderate values) and executes digit
    loops in software rather than a single ALU instruction, making heavy
    \texttt{SUM(amount)} scans roughly $5\times$ to $15\times$ more CPU-intensive
    than \texttt{SUM(amount\_minor)} on \texttt{BIGINT}. In Rust \texttt{SQLx},
    \texttt{NUMERIC} binds cleanly to \texttt{rust\_decimal::Decimal} (128-bit
    fixed decimal) or \texttt{bigdecimal::BigDecimal}.
\end{enumerate}

For single-scale ledger postings and account balances, \texttt{BIGINT} minor units
give unbeatable speed and simplicity; for FX rates, unit prices, and tax or
interest multipliers, \texttt{NUMERIC(p, s)} prevents premature rounding and
64-bit intermediate overflow.

\subsection{Time: \texttt{TIMESTAMPTZ} and the Daylight-Saving Trap}
\label{ch08:time-types}

At 01:59:59 Eastern Daylight Time on the first Sunday of November, clocks in New
York jump backward to 01:00:00 Eastern Standard Time. Suppose a ledger table uses
\texttt{TIMESTAMP WITHOUT TIME ZONE} (which SQL abbreviates as plain
\texttt{TIMESTAMP}). A customer wires funds at 01:15~EDT (\texttt{05:15 UTC}),
and fifty minutes later wires a second payment at 01:05~EST (\texttt{06:05 UTC}).
Stored as wall-clock New York \texttt{TIMESTAMP} values, the second transfer
appears to happen \emph{ten minutes before} the first: \texttt{'2026-11-01
01:05:00'} sorts ahead of \texttt{'2026-11-01 01:15:00'}. Worse, if one API node
runs with \texttt{TimeZone = 'UTC'} and an admin tool connects with
\texttt{TimeZone = 'America/Chicago'}, \texttt{TIMESTAMP WITHOUT TIME ZONE} strips
and ignores timezone offsets on input, mixing incompatible local clocks in the
same column with zero warning.

\begin{intuitionbox}
Despite its verbose SQL name, \texttt{TIMESTAMP WITH TIME ZONE}
(\texttt{TIMESTAMPTZ}) stores \textbf{zero timezone string or UTC offset on disk}.
Both \texttt{TIMESTAMP} and \texttt{TIMESTAMPTZ} occupy the exact same 8~bytes in
a heap tuple: a signed 64-bit integer counting microseconds since PostgreSQL's
epoch, \texttt{2000-01-01 00:00:00 UTC}. The sole difference happens at the input
and output boundary. When you insert \texttt{'2026-11-01 01:15:00-04'} into a
\texttt{TIMESTAMPTZ} column, PostgreSQL immediately converts that instant into UTC
microseconds and stores that single universal point on the physical timeline. On
\texttt{SELECT}, it formats those UTC microseconds into whatever session
\texttt{TimeZone} the reader requested.
\end{intuitionbox}

The rule for financial schemas has only two cases:
\begin{itemize}
  \item Every physical instant when an event occurred (\texttt{created\_at},
    \texttt{posted\_at}, \texttt{expires\_at}) must use \texttt{TIMESTAMPTZ},
    mapped in Rust \texttt{SQLx} to \texttt{DateTime<Utc>} (from \texttt{chrono})
    or \texttt{OffsetDateTime} (from \texttt{time}) with session \texttt{TimeZone}
    set to \texttt{'UTC'}.
  \item Calendar dates defined by legal contract rather than a physical instant
    (such as a wire's settlement \texttt{value\_date = '2026-11-02'} or a
    statement billing date) should use 4-byte \texttt{DATE}, never a midnight
    timestamp that shifts to the previous calendar day when viewed in another
    timezone.
\end{itemize}

\subsection{Primary Keys at Scale: Why \texttt{UUIDv4} Keys Thrash B-Trees and \texttt{UUIDv7} Flies}
\label{ch08:uuid-types}

An 8-byte \texttt{BIGINT GENERATED ALWAYS AS IDENTITY} sequence appends strictly
to the right edge of its primary-key B-tree, keeping inserts fast and compact. Yet
distributed financial systems rarely expose sequential integers as external IDs:
sequential IDs leak business volume to competitors (\texttt{entry\_id = 1042}
versus \texttt{1089} an hour later) and force Rust services or mobile clients to
make a database round-trip before knowing a new resource's identifier. That
pushes teams toward 128-bit \texttt{UUID} keys (stored natively in PostgreSQL's
16-byte \texttt{uuid} type, never as 36-byte ASCII \texttt{TEXT}).

Here is the failure mode that catches growing ledgers six months after launch. You
choose random \texttt{UUIDv4} (122 cryptographically random bits). At five million
rows, the primary-key B-tree is roughly 150~MB, sits comfortably inside an 8~GB
\texttt{shared\_buffers} pool (Chapter~2), and sustains $35{,}000$ inserts per
second. A year later, \texttt{postings} reaches two billion rows and its primary
B-tree spans 65~GB on a server with 32~GB of RAM. Insert throughput drops by an
order of magnitude while NVMe write bandwidth pins at $100\%$, even though each new
posting row is barely a hundred bytes. Three mechanisms from Chapters~2--4 combine
to cause this collapse:

\begin{enumerate}
  \item \textbf{Buffer-pool cache misses.} Every \texttt{UUIDv4} lands uniformly
    at random across the entire $[0, 2^{128})$ keyspace. When a 65~GB leaf layer
    exceeds \texttt{shared\_buffers}, nearly every single \texttt{INSERT} probes a
    cold leaf page not in memory, stalling on a random 8~KB disk read and evicting
    another page.
  \item \textbf{Midpoint page-split bloat.} As we saw in Chapter~3, random B-tree
    inserts split full leaf pages 50/50, leaving average leaf occupancy near
    $70\%$ instead of the $90\%$ fill PostgreSQL keeps when it
    detects monotonically increasing right-edge inserts.
  \item \textbf{WAL \texttt{full\_page\_writes} amplification.} Recall from
    Chapter~4 that the first modification to \emph{any} 8~KB heap or index page
    after a \texttt{CHECKPOINT} must write the entire 8~KB Full Page Image (FPI)
    into the WAL stream to survive torn writes. After a checkpoint, ten thousand
    random \texttt{UUIDv4} inserts dirty roughly ten thousand distinct B-tree leaf
    pages, turning a few hundred kilobytes of new index entries into about
    $10{,}000 \times 8\text{ KB} \approx 80\text{ MB}$ of WAL.
\end{enumerate}

RFC~9562 \textbf{\texttt{UUIDv7}} (available natively as \texttt{uuidv7()} in
PostgreSQL~18 and via \texttt{uuid::Uuid::now\_v7()} in Rust) eliminates all three
penalties without reintroducing central sequence contention. Its first 48~bits
encode the big-endian Unix timestamp in milliseconds, followed by a 4-bit version
tag (\texttt{0111}), 12~bits of sub-millisecond monotonic or random data
(\texttt{rand\_a}), a 2-bit variant (\texttt{10}), and 62~bits of cryptographic
randomness (\texttt{rand\_b}). Because PostgreSQL compares 16-byte \texttt{uuid}
values lexicographically from most-significant byte to least, every \texttt{UUIDv7}
minted during the current millisecond sorts to the \emph{far right} of the entire
historical B-tree (Figure~\ref{ch08:fig-uuidv4-vs-v7}). Concurrent Rust pods mint
collision-free IDs offline, yet all inserts converge onto a handful of hot
rightmost leaf pages pinned permanently in \texttt{shared\_buffers}.

\begin{figure}[H]
\centering
\begin{tikzpicture}[
  font=\small,
  >=Latex,
  bitbox/.style={draw=black!75, thick, minimum height=0.82cm, align=center, inner sep=2pt},
  leafcold/.style={draw=sourceblue!80!black, fill=sourceblue!8, rounded corners=2pt,
                   minimum width=1.18cm, minimum height=0.68cm, font=\scriptsize, align=center},
  leafdirty/.style={draw=trapred, fill=trapred!18, thick, rounded corners=2pt,
                    minimum width=1.18cm, minimum height=0.68cm, font=\scriptsize, align=center},
  leafhot/.style={draw=targetgreen!80!black, fill=targetgreen!20, very thick, rounded corners=2pt,
                  minimum width=1.28cm, minimum height=0.68cm, font=\scriptsize, align=center}
]
  % Top panel: UUIDv7 128-bit layout (total width = 14.0 cm)
  \node[anchor=west, font=\small\bfseries, text=structurepurple!85!black]
    at (0, 3.85) {RFC 9562 \texttt{UUIDv7} 128-Bit Big-Endian Layout (16 Bytes on Disk)};

  \node[bitbox, fill=targetgreen!18, minimum width=5.2cm, anchor=west] (ts) at (0, 3.05)
    {\textbf{48-bit Unix Epoch ms}\\ {\scriptsize monotonic B-tree sort prefix}};
  \node[bitbox, fill=structurepurple!16, minimum width=1.1cm, anchor=west] (ver) at (5.2, 3.05)
    {\textbf{4b}\\ {\scriptsize \texttt{0111}}};
  \node[bitbox, fill=sourceblue!14, minimum width=1.8cm, anchor=west] (ra) at (6.3, 3.05)
    {\textbf{12-bit}\\ {\scriptsize \texttt{rand\_a}}};
  \node[bitbox, fill=structurepurple!16, minimum width=0.9cm, anchor=west] (var) at (8.1, 3.05)
    {\textbf{2b}\\ {\scriptsize \texttt{10}}};
  \node[bitbox, fill=quotientorange!16, minimum width=5.0cm, anchor=west] (rb) at (9.0, 3.05)
    {\textbf{62-bit CSPRNG (\texttt{rand\_b})}\\ {\scriptsize collision-free across nodes}};

  \node[font=\scriptsize, text=black!70, anchor=north west] at (0, 2.58) {bit 0};
  \node[font=\scriptsize, text=black!70, anchor=north] at (5.2, 2.58) {48};
  \node[font=\scriptsize, text=black!70, anchor=north] at (6.3, 2.58) {52};
  \node[font=\scriptsize, text=black!70, anchor=north] at (8.1, 2.58) {64};
  \node[font=\scriptsize, text=black!70, anchor=north] at (9.0, 2.58) {66};
  \node[font=\scriptsize, text=black!70, anchor=north east] at (14.0, 2.58) {127};

  % Bottom-left: UUIDv4 random scatter (x: 0 .. 6.6)
  \draw[rounded corners=4pt, draw=trapred!70, fill=trapred!4] (0, -1.55) rectangle (6.65, 1.95);
  \node[font=\small\bfseries, text=trapred!85!black] at (3.32, 1.62)
    {Random \texttt{UUIDv4}: Cold Leaf Thrashing};
  \node[draw=trapred, fill=white, rounded corners=2pt, font=\scriptsize, inner sep=3pt]
    (v4in) at (3.32, 1.02) {4 concurrent inserts (uniform in $[0,2^{128})$)};

  \node[leafdirty] (l1) at (0.80, -0.15) {\texttt{0a..}\\ {\tiny 8 KB FPI}};
  \node[leafcold]  (l2) at (2.06, -0.15) {\texttt{3f..}\\ {\tiny cold}};
  \node[leafdirty] (l3) at (3.32, -0.15) {\texttt{7c..}\\ {\tiny 8 KB FPI}};
  \node[leafdirty] (l4) at (4.58, -0.15) {\texttt{b2..}\\ {\tiny 8 KB FPI}};
  \node[leafdirty] (l5) at (5.84, -0.15) {\texttt{f9..}\\ {\tiny 8 KB FPI}};

  \draw[->, thick, trapred] (v4in.south) -- (l1.north);
  \draw[->, thick, trapred] (v4in.south) -- (l3.north);
  \draw[->, thick, trapred] (v4in.south) -- (l4.north);
  \draw[->, thick, trapred] (v4in.south) -- (l5.north);

  \node[font=\scriptsize, align=center, text=black!80] at (3.32, -1.05)
    {Misses \texttt{shared\_buffers}; $\sim 50\%$ split fill;\\every touched leaf logs an 8~KB WAL FPI};

  % Bottom-right: UUIDv7 rightmost append (x: 7.35 .. 14.0)
  \draw[rounded corners=4pt, draw=targetgreen!75!black, fill=targetgreen!4] (7.35, -1.55) rectangle (14.0, 1.95);
  \node[font=\small\bfseries, text=targetgreen!75!black] at (10.68, 1.62)
    {Time-Ordered \texttt{UUIDv7}: Rightmost Append};
  \node[draw=targetgreen!80!black, fill=white, rounded corners=2pt, font=\scriptsize, inner sep=3pt]
    (v7in) at (10.68, 1.02) {4 concurrent inserts (\texttt{019a3f...} prefix)};

  \node[leafcold] (r1) at (8.12, -0.15) {\texttt{0197..}\\ {\tiny frozen}};
  \node[leafcold] (r2) at (9.38, -0.15) {\texttt{0198..}\\ {\tiny frozen}};
  \node[leafcold] (r3) at (10.64, -0.15) {\texttt{0199..}\\ {\tiny cold}};
  \node[leafhot]  (r4) at (12.45, -0.15) {\texttt{019a3f..}\\ {\tiny \textbf{Pinned RAM}}};

  \draw[->, thick, targetgreen!75!black] (v7in.south) .. controls (11.2, 0.45) .. ([xshift=-6pt]r4.north);
  \draw[->, thick, targetgreen!75!black] (v7in.south) .. controls (11.6, 0.45) .. ([xshift=-2pt]r4.north);
  \draw[->, thick, targetgreen!75!black] (v7in.south) .. controls (12.0, 0.45) .. ([xshift=2pt]r4.north);
  \draw[->, thick, targetgreen!75!black] (v7in.south) .. controls (12.4, 0.45) .. ([xshift=6pt]r4.north);

  \node[font=\scriptsize, align=center, text=black!80] at (10.68, -1.05)
    {$>99\%$ buffer hit rate on rightmost leaf;\\$\sim 90\%+$ fill factor; 1 FPI per checkpoint};
\end{tikzpicture}
\caption{Bit layout of RFC~9562 \texttt{UUIDv7} (top) and its effect on a multi-gigabyte B-tree (bottom). Random \texttt{UUIDv4} keys scatter writes across cold disk pages and amplify post-checkpoint WAL full-page images, whereas \texttt{UUIDv7} funnels concurrent inserts onto the hot rightmost leaf in \texttt{shared\_buffers}.}
\label{ch08:fig-uuidv4-vs-v7}
\end{figure}

\section{Encoding Domain Invariants Directly in the Schema}
\label{ch08:invariants}

``Our Rust service validates every request struct before calling SQL, so the
database only needs plain columns.'' That sentence precedes nearly every silent
data-corruption postmortem. Application code is only one doorway into a database:
settlement batch jobs, emergency \texttt{psql} backfills, schema migrations, and
concurrent network retries all write to the same tables. Worse, under \texttt{READ
COMMITTED} isolation (Chapter~6), two Rust handlers that validate against separate
snapshots can each see valid preconditions and still commit a combined state that
violates business law. Pushing invariants into PostgreSQL turns application bugs
from silent ledger corruption into loud, rolled-back \texttt{ERROR} responses.

\subsection{\texttt{NOT NULL}, \texttt{CHECK}, Domains, and Atomic Idempotency Keys}
\label{ch08:scalar-invariants}

Recall from Chapter~1 how three-valued logic treats \texttt{NULL} inside a
\texttt{CHECK} constraint: a row passes \texttt{CHECK (expr)} whenever
\texttt{expr} evaluates to \texttt{TRUE} \emph{or} \texttt{UNKNOWN}. If you write
\texttt{amount\_minor BIGINT CHECK (amount\_minor > 0)} without \texttt{NOT NULL},
inserting \texttt{NULL} evaluates \texttt{NULL > 0} to \texttt{UNKNOWN} and lands
cleanly on the heap page. Always pair \texttt{NOT NULL} with \texttt{CHECK}.

For constrained text categories, choose deliberately between three tools:
\begin{itemize}
  \item \textbf{\texttt{CREATE TYPE ... AS ENUM}} stores each label as a compact
    4-byte OID on disk and excels for closed state machines whose values are tied
    to code branches (\texttt{'pending'}, \texttt{'posted'}, \texttt{'voided'}).
    However, PostgreSQL enums cannot drop obsolete labels without recreating the
    type.
  \item \textbf{\texttt{DOMAIN}s or \texttt{CHECK}-constrained \texttt{TEXT}}
    (such as \texttt{CHAR(3)} checked against \texttt{'\textasciicircum{}[A-Z]\{3\}\$'})
    validate format uniformly across dozens of tables while allowing new currencies
    without DDL catalog locks. When currencies carry per-asset metadata such as
    minor-unit scale ($2$ for \texttt{USD}, $0$ for \texttt{JPY}), a foreign key to
    a small \texttt{assets} table (\texttt{currency}, \texttt{minor\_scale}) stays
    permanently cached in \texttt{shared\_buffers}.
  \item \textbf{Storage-enforced idempotency keys.} When a mobile app or upstream
    gateway sends \texttt{POST /v1/transfers} and its TCP socket times out while
    waiting for the HTTP response, the client cannot tell whether the commit
    succeeded; it \emph{must} retry with the same client-generated
    \texttt{idempotency\_key}. Checking \texttt{SELECT 1 FROM journal\_entries
    WHERE client\_id = \$1 AND idempotency\_key = \$2} in Rust before inserting
    suffers a classic Time-of-Check to Time-of-Use race if two retries arrive
    simultaneously. Declaring \texttt{UNIQUE (client\_id, idempotency\_key)} makes
    PostgreSQL's B-tree uniqueness machinery (\texttt{\_bt\_check\_unique}) inspect
    even uncommitted in-flight index tuples, block the second retry on the first
    transaction's lock, and resolve atomically via \texttt{ON CONFLICT DO NOTHING}.
\end{itemize}

\subsection{Temporal Non-Overlapping Invariants with \texttt{EXCLUDE USING gist}}
\label{ch08:gist-exclude}

Suppose each savings or margin account has a tiered interest-rate schedule valid
over a time interval. Two administrative jobs concurrently insert schedules for
account~$A$: one sets $4.50\%$ over $[\text{Jun~1}, \text{Dec~1})$ and another sets
$5.00\%$ over $[\text{Sep~1}, \text{Mar~1})$. A nightly accrual query joins
\texttt{accounts} to \texttt{account\_rate\_schedules WHERE now() <@
effective\_during}, matches \emph{both} rows throughout autumn, and pays $9.50\%$
interest. No \texttt{UNIQUE} constraint on \texttt{(account\_id, effective\_from)}
can catch partial interval overlap because the start timestamps differ.

Chapter~3 introduced Generalized Search Trees (GiST) as balanced trees over
arbitrary bounding predicates rather than total orders. Paired with the bundled
\texttt{btree\_gist} extension (which teaches GiST to test equality \texttt{=} on
ordinary scalar types alongside range overlap \texttt{\&\&}), PostgreSQL's
\textbf{exclusion constraint} rejects any two rows where \emph{all} listed
operator comparisons return true simultaneously:

\begin{lstlisting}[language=SQL]
CREATE EXTENSION IF NOT EXISTS btree_gist;

CREATE TABLE account_rate_schedules (
    schedule_id      UUID PRIMARY KEY DEFAULT uuidv7(),
    account_id       UUID NOT NULL,
    annual_bps       INTEGER NOT NULL CHECK (annual_bps BETWEEN 0 AND 10000),
    effective_during TSTZRANGE NOT NULL,
    CONSTRAINT non_empty_window CHECK (NOT isempty(effective_during)),
    CONSTRAINT no_overlapping_rates EXCLUDE USING gist (
        account_id       WITH =,
        effective_during WITH &&
    )
);
\end{lstlisting}

Notice the range type \texttt{TSTZRANGE} and PostgreSQL's canonical half-open
convention \texttt{'[)'} (closed lower bound, open upper bound). When an old rate
schedule \texttt{[2026-01-01Z, 2026-06-01Z)} (constructed via \texttt{tstzrange})
ends at midnight June~1 and its replacement \texttt{[2026-06-01Z, infinity)}
begins at that exact microsecond, the two intervals touch with zero gap yet share
zero instants, so \texttt{\&\&} evaluates to \texttt{false}. Even under
\texttt{READ COMMITTED}, two concurrent transactions inserting overlapping windows
for the same \texttt{account\_id} collide inside the GiST index just like a
duplicate primary key: one commits and the other aborts with SQLSTATE
\texttt{23P01} (\texttt{exclusion\_violation}).

\section{A Double-Entry Financial Ledger in PostgreSQL and Rust}
\label{ch08:ledger}

Beginners often model money with a single mutable column: \texttt{UPDATE accounts
SET balance = balance - 5000 WHERE id = \$1} followed by \texttt{UPDATE accounts
SET balance = balance + 5000 WHERE id = \$2}. Three production hazards make that
design unacceptable for a real ledger:
\begin{enumerate}
  \item \textbf{Destroyed history.} Overwriting \texttt{balance} in place erases
    \emph{why} the balance changed, making regulatory audits, dispute reversals,
    and point-in-time historical reconstructions impossible once \texttt{VACUUM}
    reclaims old tuple versions (Chapter~5).
  \item \textbf{Unchecked money creation.} Nothing in the schema ties the two
    \texttt{UPDATE} statements together. A bug in a new promo-credit handler that
    runs one \texttt{UPDATE} without its counterparty silently alters the total
    supply of customer funds.
  \item \textbf{Hot-row contention without audit legs.} Separating immutable
    event recording from balance caching lets us enforce conservation of money at
    the database kernel level while keeping point balance reads $O(1)$.
\end{enumerate}

\subsection{Conservation of Value and \texttt{DEFERRABLE INITIALLY DEFERRED} Triggers}
\label{ch08:deferred-triggers}

Luca Pacioli codified the cure in Venice in 1494: \textbf{double-entry
bookkeeping}. Every financial event is recorded as an immutable
\texttt{journal\_entry} header accompanied by two or more immutable
\texttt{postings} (legs) in the same currency whose signed minor-unit amounts sum
to zero:
\[
  \forall\, e \in \text{journal\_entries}:\qquad
  \abs{P(e)} \ge 2
  \quad\text{and}\quad
  \sum_{p \in P(e)} p.\texttt{signed\_amount\_minor} = 0.
\]
When Alice sends $\$50.00$ ($5{,}000$ cents) to Bob, we do not invent $5{,}000$
cents in Bob's row; we append one leg of $-5000$ against Alice's liability account
and one leg of $+5000$ against Bob's liability account under the same
\texttt{entry\_id}. If a $1\%$ platform fee applies, the same entry has three legs:
$-5000$ (Alice), $+4950$ (Bob), and $+50$ (Platform Fee Revenue), summing cleanly
to $-5000 + 4950 + 50 = 0$.

How do we force PostgreSQL itself to reject any transaction that violates this
zero-sum balance law? A table-level \texttt{CHECK} constraint can inspect only the
columns of a \emph{single} row. And an ordinary \texttt{AFTER INSERT} trigger fails
immediately: the moment you insert Alice's $-5000$ leg (before inserting Bob's
$+5000$ leg), the running sum for \texttt{entry\_id} is $-5000 \ne 0$!

PostgreSQL solves multi-row invariants across statement boundaries with a
\textbf{\texttt{CONSTRAINT TRIGGER}} declared \textbf{\texttt{DEFERRABLE INITIALLY
DEFERRED}}. Instead of firing mid-transaction when only half the legs exist,
PostgreSQL queues a lightweight trigger event in backend memory and executes the
verification function \emph{when the client issues \texttt{COMMIT}}, right before
flushing the commit record to WAL (Figure~\ref{ch08:fig-gist-and-deferred}). Let
us write out the complete ledger schema, including an immutability trigger on
\texttt{postings} and a materialized \texttt{cached\_balance\_minor} on
\texttt{accounts} protected against negative overdrafts:

\begin{lstlisting}[language=SQL]
CREATE TABLE accounts (
    account_id           UUID PRIMARY KEY DEFAULT uuidv7(),
    currency             CHAR(3) NOT NULL CHECK (currency ~ '^[A-Z]{3}$'),
    cached_balance_minor BIGINT NOT NULL DEFAULT 0,
    allow_overdraft      BOOLEAN NOT NULL DEFAULT FALSE,
    created_at           TIMESTAMPTZ NOT NULL DEFAULT now(),
    CONSTRAINT non_negative_unless_overdraft
        CHECK (allow_overdraft OR cached_balance_minor >= 0)
);

CREATE TABLE journal_entries (
    entry_id        UUID PRIMARY KEY DEFAULT uuidv7(),
    client_id       UUID NOT NULL,
    idempotency_key TEXT NOT NULL CHECK (length(idempotency_key) <= 128),
    currency        CHAR(3) NOT NULL CHECK (currency ~ '^[A-Z]{3}$'),
    description     TEXT NOT NULL,
    created_at      TIMESTAMPTZ NOT NULL DEFAULT now(),
    CONSTRAINT uq_client_idempotency UNIQUE (client_id, idempotency_key)
);

CREATE TABLE postings (
    posting_id          UUID PRIMARY KEY DEFAULT uuidv7(),
    entry_id            UUID NOT NULL REFERENCES journal_entries(entry_id),
    account_id          UUID NOT NULL REFERENCES accounts(account_id),
    signed_amount_minor BIGINT NOT NULL CHECK (signed_amount_minor <> 0),
    created_at          TIMESTAMPTZ NOT NULL DEFAULT now()
);

CREATE INDEX idx_postings_entry   ON postings (entry_id);
CREATE INDEX idx_postings_account ON postings (account_id, created_at DESC);

CREATE OR REPLACE FUNCTION forbid_posting_mutation()
RETURNS TRIGGER LANGUAGE plpgsql AS $$
BEGIN
    RAISE EXCEPTION 'postings are append-only; post a reversing entry';
END;
$$;

CREATE TRIGGER trg_postings_immutable
BEFORE UPDATE OR DELETE ON postings
FOR EACH ROW EXECUTE FUNCTION forbid_posting_mutation();

CREATE OR REPLACE FUNCTION verify_entry_zero_sum()
RETURNS TRIGGER LANGUAGE plpgsql AS $$
DECLARE
    v_entry_id UUID := COALESCE(NEW.entry_id, OLD.entry_id);
    v_legs     BIGINT;
    v_net_sum  BIGINT;
    v_bad_ccy  BIGINT;
BEGIN
    SELECT COUNT(*), COALESCE(SUM(p.signed_amount_minor), 0),
           COUNT(*) FILTER (WHERE a.currency <> je.currency)
      INTO v_legs, v_net_sum, v_bad_ccy
      FROM journal_entries je
      LEFT JOIN postings p ON p.entry_id = je.entry_id
      LEFT JOIN accounts a ON a.account_id = p.account_id
     WHERE je.entry_id = v_entry_id;

    IF v_legs < 2 OR v_net_sum <> 0 OR v_bad_ccy > 0 THEN
        RAISE EXCEPTION USING
            ERRCODE = '23514',
            MESSAGE = 'unbalanced entry ' || v_entry_id
                   || ': legs=' || v_legs || ', sum=' || v_net_sum
                   || ', ccy_mismatch=' || v_bad_ccy;
    END IF;
    RETURN NULL;
END;
$$;

CREATE CONSTRAINT TRIGGER trg_verify_entry_zero_sum
AFTER INSERT OR UPDATE OR DELETE ON postings
DEFERRABLE INITIALLY DEFERRED
FOR EACH ROW EXECUTE FUNCTION verify_entry_zero_sum();
\end{lstlisting}

\begin{trapbox}
\texttt{verify\_entry\_zero\_sum} checks \emph{three} things
together: \texttt{v\_legs < 2}, \texttt{v\_net\_sum <> 0}, and
\texttt{v\_bad\_ccy > 0}. Without \texttt{LEFT JOIN} and \texttt{v\_legs < 2}, an
entry whose postings were somehow deleted (if someone disabled the immutability
trigger during maintenance) would return \texttt{SUM(...) = NULL}, which naive
\texttt{IF v\_net\_sum <> 0} evaluates to \texttt{UNKNOWN} and lets slip past. And
without checking \texttt{a.currency <> je.currency}, a buggy handler could debit
5000~JPY from a yen account and credit 5000~USD cents to a dollar account while
still summing to zero numerically. Finally, attach a second deferred check on
\texttt{journal\_entries} (or always insert postings in the same helper) so an
empty header with zero postings cannot commit alone.
\end{trapbox}

\begin{figure}[H]
\centering
\begin{tikzpicture}[
  font=\small,
  >=Latex,
  stepbox/.style={draw=sourceblue!80!black, fill=sourceblue!8, rounded corners=2pt,
                  minimum height=0.68cm, inner sep=3pt, font=\scriptsize, align=center}
]
  % Panel A: GiST Exclusion Geometry (top, y: 2.3 .. 4.9)
  \draw[rounded corners=4pt, draw=structurepurple!75, fill=structurepurple!4]
    (0, 2.25) rectangle (14.2, 4.95);
  \node[font=\small\bfseries, text=structurepurple!85!black, anchor=west]
    at (0.3, 4.62) {A. \texttt{EXCLUDE USING gist (account\_id WITH =, effective\_during WITH \&\&)}};

  % Time axis
  \draw[->, thick, black!75] (1.2, 2.70) -- (13.6, 2.70) node[right, font=\scriptsize] {$t$ (UTC)};
  \foreach \x/\lbl in {2.0/Jan 1, 5.2/Apr 1, 7.4/Jun 1, 10.2/Sep 1, 12.6/Dec 1} {
    \draw[black!70] (\x, 2.62) -- (\x, 2.78);
    \node[font=\scriptsize, below] at (\x, 2.62) {\lbl};
  }

  % Valid half-open intervals [Jan 1, Jun 1) and [Jun 1, Dec 1)
  \draw[very thick, targetgreen!80!black] (2.0, 3.95) -- (7.32, 3.95);
  \fill[targetgreen!80!black] (2.0, 3.95) circle (2.3pt);
  \draw[thick, targetgreen!80!black, fill=white] (7.4, 3.95) circle (2.3pt);
  \node[font=\scriptsize, text=targetgreen!75!black, above] at (4.65, 3.98)
    {Row 1: \texttt{[Jan 1, Jun 1)} @ 4.50\%};

  \draw[very thick, sourceblue!85!black] (7.48, 3.95) -- (12.52, 3.95);
  \fill[sourceblue!85!black] (7.4, 3.95) circle (2.3pt);
  \draw[thick, sourceblue!85!black, fill=white] (12.6, 3.95) circle (2.3pt);
  \node[font=\scriptsize, text=sourceblue!85!black, above] at (10.1, 3.98)
    {Row 2: \texttt{[Jun 1, Dec 1)} @ 4.75\% (no \texttt{\&\&} at Jun 1)};

  % Rejected overlapping interval [Apr 1, Sep 1)
  \draw[very thick, trapred, dashed] (5.2, 3.25) -- (10.2, 3.25);
  \fill[trapred] (5.2, 3.25) circle (2.3pt);
  \draw[thick, trapred, fill=white] (10.2, 3.25) circle (2.3pt);
  \fill[trapred!20, opacity=0.55] (5.2, 3.15) rectangle (10.2, 3.35);
  \node[font=\scriptsize\bfseries, text=trapred, right] at (10.35, 3.25)
    {Rejected (\texttt{23P01}): overlaps!};

  % Panel B: Deferred Constraint Trigger Timeline (bottom, y: -1.15 .. 1.95)
  \draw[rounded corners=4pt, draw=sourceblue!75, fill=sourceblue!3]
    (0, -1.15) rectangle (14.2, 1.95);
  \node[font=\small\bfseries, text=sourceblue!85!black, anchor=west]
    at (0.3, 1.62) {B. \texttt{DEFERRABLE INITIALLY DEFERRED} Trigger Timeline Across a Transfer};

  \node[stepbox, minimum width=1.55cm] (s1) at (1.15, 0.70) {\texttt{BEGIN}\\ {\tiny $\Sigma = 0$}};
  \node[stepbox, minimum width=2.35cm, draw=quotientorange!90!black, fill=quotientorange!12]
    (s2) at (3.75, 0.70) {Leg 1: $-5000$\\ {\tiny $\Sigma = -5000$ (queued)}};
  \node[stepbox, minimum width=2.35cm]
    (s3) at (6.75, 0.70) {Leg 2: $+5000$\\ {\tiny $\Sigma = 0$ (queued)}};
  \node[stepbox, minimum width=1.85cm, draw=structurepurple, fill=structurepurple!14]
    (s4) at (9.50, 0.70) {\texttt{COMMIT}\\ {\tiny fire triggers}};
  \node[stepbox, minimum width=2.85cm, draw=targetgreen!80!black, fill=targetgreen!16]
    (s5) at (12.48, 0.70) {$\Sigma = 0 \land \text{legs} \ge 2$\\ {\tiny flush WAL + ack}};

  \draw[->, thick] (s1) -- (s2);
  \draw[->, thick] (s2) -- (s3);
  \draw[->, thick] (s3) -- (s4);
  \draw[->, thick] (s4) -- (s5);

  \node[font=\scriptsize, text=black!75, align=center] at (7.1, -0.45)
    {Mid-transaction imbalance after Leg~1 is permitted; if the legs do not sum to zero at \texttt{COMMIT},\\PostgreSQL aborts the transaction before writing any commit record to WAL.};
\end{tikzpicture}
\caption{Two declarative PostgreSQL invariants that replace fragile application-side checks: (A) half-open \texttt{TSTZRANGE} intervals enforced by \texttt{EXCLUDE USING gist}, and (B) commit-time conservation verification via a deferred constraint trigger.}
\label{ch08:fig-gist-and-deferred}
\end{figure}

\subsection{Deadlock-Free Transfers in Rust (\texttt{Axum} + \texttt{SQLx})}
\label{ch08:rust-ledger}

Now let us wire our schema to a Rust \texttt{Axum} service using \texttt{SQLx}. Two
subtle concurrency failures await careless transfer code even when the schema is
correct:

\begin{enumerate}
  \item \textbf{Opposite-order deadlocks (Chapter~6).} Suppose Task~1 transfers
    $\$10$ from Alice to Bob while Task~2 simultaneously transfers $\$5$ from Bob
    to Alice. If each handler locks \texttt{from\_account} first and
    \texttt{to\_account} second, Task~1 holds Alice's row lock and waits on Bob,
    while Task~2 holds Bob's row lock and waits on Alice until PostgreSQL's
    \texttt{deadlock\_timeout} kills one transaction. Sorting the account UUIDs in
    canonical ascending order (\texttt{ORDER BY account\_id}) before acquiring
    locks makes deadlocks across account pairs mathematically impossible, because a
    strict total order has no cycles.
  \item \textbf{Over-strong row locks blocking foreign keys.} Recall from Chapter~6
    that \texttt{INSERT INTO postings} automatically acquires a \texttt{FOR KEY
    SHARE} tuple lock on the referenced \texttt{accounts} row to make sure nobody
    deletes that account mid-transaction. If you lock the account rows with
    \texttt{SELECT ... FOR UPDATE}, PostgreSQL assumes you might change the
    primary key \texttt{account\_id} itself, which conflicts with \texttt{FOR KEY
    SHARE}. Instead, use \textbf{\texttt{FOR NO KEY UPDATE}}: it still serializes
    concurrent \texttt{UPDATE accounts SET cached\_balance\_minor = ...} mutations
    on the same account, while remaining compatible with \texttt{FOR KEY SHARE}
    foreign-key checks.
\end{enumerate}

Listing~\ref{ch08:lst-rust-transfer} puts every piece together: input validation,
atomic idempotency deduplication, canonical \texttt{FOR NO KEY UPDATE} lock
ordering, double-entry leg insertion, cached balance updates, and deferred-trigger
verification on \texttt{tx.commit().await}.

\begin{lstlisting}[language=Rust, caption={Idempotent, deadlock-free double-entry transfer in Rust with \texttt{SQLx}.}, label={ch08:lst-rust-transfer}]
use sqlx::{PgPool, Postgres, Transaction};
use uuid::Uuid;

#[derive(Debug)]
pub struct TransferCmd {
    pub client_id: Uuid,
    pub idempotency_key: String,
    pub from_account: Uuid,
    pub to_account: Uuid,
    pub currency: String,
    pub amount_minor: i64,
    pub description: String,
}

#[derive(Debug, thiserror::Error)]
pub enum LedgerError {
    #[error("invalid transfer parameters")]
    InvalidInput,
    #[error("account missing or currency mismatch")]
    AccountMismatch,
    #[error("insufficient funds")]
    InsufficientFunds,
    #[error("database error: {0}")]
    Db(#[from] sqlx::Error),
}

#[derive(sqlx::FromRow)]
struct LockedAccount {
    account_id: Uuid,
    currency: String,
    cached_balance_minor: i64,
    allow_overdraft: bool,
}

pub async fn execute_transfer(
    pool: &PgPool,
    cmd: &TransferCmd,
) -> Result<Uuid, LedgerError> {
    if cmd.amount_minor <= 0 || cmd.from_account == cmd.to_account {
        return Err(LedgerError::InvalidInput);
    }

    let mut tx: Transaction<'_, Postgres> = pool.begin().await?;

    // 1. Atomic idempotency claim via UNIQUE (client_id, idempotency_key).
    let claimed: Option<Uuid> = sqlx::query_scalar(
        "INSERT INTO journal_entries \
            (client_id, idempotency_key, currency, description) \
         VALUES ($1, $2, $3, $4) \
         ON CONFLICT (client_id, idempotency_key) DO NOTHING \
         RETURNING entry_id",
    )
    .bind(cmd.client_id)
    .bind(&cmd.idempotency_key)
    .bind(&cmd.currency)
    .bind(&cmd.description)
    .fetch_optional(&mut *tx)
    .await?;

    let Some(entry_id) = claimed else {
        // Retry detected: return the committed entry_id cleanly.
        let existing: Uuid = sqlx::query_scalar(
            "SELECT entry_id FROM journal_entries \
             WHERE client_id = $1 AND idempotency_key = $2",
        )
        .bind(cmd.client_id)
        .bind(&cmd.idempotency_key)
        .fetch_one(&mut *tx)
        .await?;
        tx.rollback().await?;
        return Ok(existing);
    };

    // 2. Canonical lock order prevents A->B vs B->A deadlocks.
    let mut ids = [cmd.from_account, cmd.to_account];
    ids.sort_unstable();

    let locked: Vec<LockedAccount> = sqlx::query_as(
        "SELECT account_id, currency, cached_balance_minor, \
                allow_overdraft \
         FROM accounts \
         WHERE account_id = ANY($1) \
         ORDER BY account_id \
         FOR NO KEY UPDATE",
    )
    .bind(&ids[..])
    .fetch_all(&mut *tx)
    .await?;

    if locked.len() != 2
        || locked.iter().any(|a| a.currency.trim() != cmd.currency)
    {
        return Err(LedgerError::AccountMismatch);
    }

    let sender = locked
        .iter()
        .find(|a| a.account_id == cmd.from_account)
        .ok_or(LedgerError::AccountMismatch)?;

    if !sender.allow_overdraft
        && sender.cached_balance_minor < cmd.amount_minor
    {
        return Err(LedgerError::InsufficientFunds);
    }

    // 3. Append both immutable postings in a single multi-row INSERT.
    sqlx::query(
        "INSERT INTO postings \
            (entry_id, account_id, signed_amount_minor) \
         VALUES ($1, $2, $3), ($1, $4, $5)",
    )
    .bind(entry_id)
    .bind(cmd.from_account)
    .bind(-cmd.amount_minor)
    .bind(cmd.to_account)
    .bind(cmd.amount_minor)
    .execute(&mut *tx)
    .await?;

    // 4. Update cached balances while holding FOR NO KEY UPDATE locks.
    sqlx::query(
        "UPDATE accounts \
         SET cached_balance_minor = cached_balance_minor + CASE \
             WHEN account_id = $1 THEN -$3 \
             WHEN account_id = $2 THEN  $3 \
         END \
         WHERE account_id IN ($1, $2)",
    )
    .bind(cmd.from_account)
    .bind(cmd.to_account)
    .bind(cmd.amount_minor)
    .execute(&mut *tx)
    .await?;

    // 5. Commit fires trg_verify_entry_zero_sum inside PostgreSQL!
    tx.commit().await?;
    Ok(entry_id)
}
\end{lstlisting}

\section{Declarative Partitioning for Billion-Row Ledgers}
\label{ch08:partitioning}

Imagine our ledger sustains $150$ transfers ($300$ postings) per second around the
clock: $300 \times 86{,}400 \times 365 \approx 9.5$ billion new \texttt{postings}
rows a year. At roughly 100 bytes per row (tuple header, line pointer, and five
columns), that is close to a terabyte of heap pages a year before any index.
Keeping years of that inside a single physical relation causes three operational
failures:

\begin{enumerate}
  \item \textbf{Cold secondary indexes.} Even with \texttt{UUIDv7} keeping primary
    key inserts localized, secondary indexes such as \texttt{(account\_id,
    created\_at DESC)} span years of history across hundreds of gigabytes, evicting
    hot pages from \texttt{shared\_buffers}.
  \item \textbf{Multi-hour \texttt{VACUUM} freeze passes.} As we learned in
    Chapter~5, even an append-only table that never runs \texttt{UPDATE} or
    \texttt{DELETE} still requires periodic \texttt{autovacuum} freeze scans to
    mark heap pages all-frozen in the Visibility Map before the 32-bit transaction
    ID horizon wraps around. Scanning a half-terabyte single heap file saturates
    storage I/O for hours.
  \item \textbf{Retention archival via \texttt{DELETE} bloats disk.} When
    compliance policy moves postings older than seven years to cold object storage,
    running \texttt{DELETE FROM postings WHERE created\_at < '2019-11-01'} on
    $250\text{ million}$ rows does \emph{not} return bytes to the filesystem: it
    marks $250\text{ million}$ dead tuples, writes tens of gigabytes of WAL, and
    leaves fragmented pages behind.
\end{enumerate}

PostgreSQL's \textbf{declarative range partitioning} solves all three problems by
presenting one logical parent table (\texttt{postings}) backed by independent
monthly physical child tables (\texttt{postings\_2026\_10},
\texttt{postings\_2026\_11}, and so on), each owning its own heap files, Visibility
Map, and local B-tree indexes:

\begin{lstlisting}[language=SQL]
CREATE TABLE postings (
    posting_id          UUID NOT NULL DEFAULT uuidv7(),
    entry_id            UUID NOT NULL,
    account_id          UUID NOT NULL,
    signed_amount_minor BIGINT NOT NULL CHECK (signed_amount_minor <> 0),
    created_at          TIMESTAMPTZ NOT NULL DEFAULT now(),
    PRIMARY KEY (posting_id, created_at)
) PARTITION BY RANGE (created_at);

CREATE TABLE postings_2026_11 PARTITION OF postings
    FOR VALUES FROM ('2026-11-01 00:00:00+00')
                 TO ('2026-12-01 00:00:00+00');

CREATE TABLE postings_2026_12 PARTITION OF postings
    FOR VALUES FROM ('2026-12-01 00:00:00+00')
                 TO ('2027-01-01 00:00:00+00');

CREATE INDEX idx_postings_acct_time
    ON postings (account_id, created_at DESC);
\end{lstlisting}

Look at what changes inside the engine:
\begin{itemize}
  \item \textbf{Plan-time and execution-time partition pruning.} When a customer
    loads their November statement (\texttt{WHERE account\_id = \$1} with November
    \texttt{created\_at} bounds), the query planner (Chapter~7) compares those
    predicates against each child's \texttt{FOR VALUES} range and prunes $83$ out
    of $84$ monthly partitions before touching a single index page. Even when bounds
    come from runtime parameters or join keys, PostgreSQL's \texttt{Append} executor
    prunes non-matching partitions dynamically at execution startup. Only the
    current month's B-trees need to stay warm in \texttt{shared\_buffers}.
  \item \textbf{Localized \texttt{autovacuum}.} Once November ends, all pages in
    \texttt{postings\_2026\_11} are frozen once in the Visibility Map and never
    scanned again by \texttt{VACUUM}.
  \item \textbf{Instantaneous zero-bloat archival.} Expiring the oldest month
    never runs \texttt{DELETE}:
\begin{lstlisting}[language=SQL]
ALTER TABLE postings DETACH PARTITION postings_2019_10 CONCURRENTLY;
-- Export standalone table postings_2019_10 to Parquet/archive, then:
DROP TABLE postings_2019_10;
\end{lstlisting}
    \texttt{DETACH PARTITION CONCURRENTLY} updates system catalogs without blocking
    concurrent reads or inserts on live months, and \texttt{DROP TABLE} unlinks the
    underlying 1~GB OS segment files in milliseconds with zero dead tuples.
\end{itemize}

\begin{trapbox}
Notice that our partitioned table declared \texttt{PRIMARY KEY (posting\_id,
created\_at)} rather than \texttt{PRIMARY KEY (posting\_id)}. PostgreSQL rejects
any \texttt{UNIQUE} or \texttt{PRIMARY KEY} constraint on a partitioned table
unless \emph{every column of the partition key appears inside the constraint}.
From Chapter~3 you can see why: each monthly partition owns a separate local
B-tree, and enforcing global uniqueness on \texttt{posting\_id} alone would force
every single \texttt{INSERT} to probe all $84$ monthly B-trees. Keep
\texttt{journal\_entries} unpartitioned (or partitioned with \texttt{created\_at}
included in the idempotency window) if you need a single global B-tree on
\texttt{(client\_id, idempotency\_key)}.
\end{trapbox}

\section{High Availability: Physical Streaming and Logical Replication}
\label{ch08:replication}

A RAID array or cloud block volume protects one server against a single drive
failure, not against a kernel panic, motherboard fault, or availability-zone power
loss. Yet simply standing up a standby server is not enough: depending on two
configuration parameters, an automatic failover can still silently erase the last
few milliseconds of committed customer deposits, or a read-balanced API endpoint
can report a stale balance right after a transfer succeeds.

\subsection{Physical Streaming WAL Replication and Standby Conflicts}
\label{ch08:physical-replication}

Recall from Chapter~4 that PostgreSQL's Write-Ahead Log records physical changes
keyed by relation file node, 8~KB block number, and page offset, ordered globally
by a monotonically increasing 64-bit Log Sequence Number (\texttt{XLogRecPtr}, or
LSN). \textbf{Physical streaming replication} turns that crash-recovery log into a
continuous network stream:

\begin{enumerate}
  \item On the primary, a \texttt{walsender} backend process reads freshly
    generated WAL bytes straight from WAL buffers or \texttt{pg\_wal} and streams
    them over TCP to each connected standby. Attaching a \textbf{physical
    replication slot} records how far each standby has consumed so the primary
    never recycles a 16~MB WAL segment a disconnected replica still needs (always
    alert on \texttt{pg\_replication\_slots} lag: an orphaned slot on a dead
    standby will retain WAL indefinitely until the primary's disk fills $100\%$).
  \item On the standby, a \texttt{walreceiver} process writes incoming WAL bytes to
    the standby's own \texttt{pg\_wal} directory and \texttt{fsync}s them, while
    the standby's \texttt{startup} process continuously replays those records onto
    the standby's 8~KB heap and index pages.
\end{enumerate}

Because physical replication ships exact page modifications, a physical standby is
a \textbf{byte-for-byte clone} of the primary: identical 8~KB slotted pages,
identical \texttt{xmin}/\texttt{xmax} transaction IDs, and identical B-tree leaves
across the entire cluster. With \texttt{hot\_standby = on}, read-only queries can
execute directly against the standby while replay runs underneath them.

That concurrency creates a subtle clash with Chapter~5's \texttt{VACUUM}. Suppose a
reporting query runs for forty seconds on the hot standby using snapshot
\texttt{xmin = 50000}. Meanwhile on the primary, no old query is running, so
\texttt{VACUUM} reclaims dead tuples from transaction \texttt{50010} and emits a
\texttt{HEAP2\_PRUNE} WAL record. When the standby's \texttt{startup} process
reaches that WAL record, replaying it immediately would delete tuples still visible
to the standby's running query. PostgreSQL resolves this \textbf{recovery conflict}
via \texttt{max\_standby\_streaming\_delay} (default \texttt{30s}): the standby
pauses WAL replay up to that threshold to let the read query finish, and if the
query is still running when the timer expires, cancels the read query with
\texttt{ERROR: canceling statement due to conflict with recovery}. Enabling
\texttt{hot\_standby\_feedback = on} tells the standby to advertise its oldest
active \texttt{xmin} back to the primary's \texttt{ProcArray}, preventing query
cancellations at the cost of holding back \texttt{VACUUM} cleanup on the primary
whenever a replica runs multi-minute analytics.

\subsection{The \texttt{synchronous\_commit} Durability and Visibility Ladder}
\label{ch08:sync-commit-ladder}

When a Rust handler calls \texttt{tx.commit().await}, how far through the
replication pipeline must the commit LSN travel before the primary returns
\texttt{CommandComplete} to the client? PostgreSQL lets you choose the exact point
on the latency-versus-durability frontier, either globally in
\texttt{postgresql.conf} or for one transaction at a time with
\texttt{SET LOCAL}. Two settings work as a pair: \texttt{synchronous\_commit}
says how far the commit must travel, and \texttt{synchronous\_standby\_names}
says which standbys count (for example,
\texttt{ANY 1 (az2\_standby, az3\_standby)} across three availability zones).

Each standby continuously sends feedback messages reporting three LSN watermarks,
all visible in the \texttt{pg\_stat\_replication} view:
\begin{itemize}
  \item \texttt{write\_lsn}: WAL copied into the standby's OS page cache;
  \item \texttt{flush\_lsn}: WAL \texttt{fsync}ed to the standby's NVMe drive;
  \item \texttt{replay\_lsn}: WAL applied to the standby's buffer pages, and
  therefore visible to standby snapshots.
\end{itemize}
Figure~\ref{ch08:fig-replication-ladder} and
Table~\ref{ch08:tab-sync-commit} compare all five levels.

\begin{figure}[H]
\centering
\begin{tikzpicture}[
  font=\small,
  >=Latex,
  stage/.style={draw=black!75, rounded corners=3pt, minimum width=2.30cm,
                minimum height=1.05cm, align=center, font=\scriptsize, inner sep=3pt},
  rung/.style={font=\scriptsize\bfseries, anchor=west}
]
  % Primary box (left) and Standby box (right), total width = 14.2 cm
  \draw[rounded corners=4pt, draw=sourceblue!80!black, fill=sourceblue!5]
    (0, 1.35) rectangle (5.55, 3.35);
  \node[font=\small\bfseries, text=sourceblue!85!black] at (2.77, 3.05)
    {Primary Node (AZ-1)};

  \node[stage, fill=white] (pwal) at (1.40, 2.12)
    {1. WAL Buffer\\ {\tiny \texttt{off} ($\sim 0.05\text{ ms}$)}};
  \node[stage, fill=sourceblue!16] (pdisk) at (4.15, 2.12)
    {2. Local NVMe\\ {\tiny \texttt{local} ($\sim 0.3\text{ ms}$)}};
  \draw[->, thick] (pwal) -- (pdisk);

  \draw[rounded corners=4pt, draw=targetgreen!80!black, fill=targetgreen!5]
    (6.35, 1.35) rectangle (14.20, 3.35);
  \node[font=\small\bfseries, text=targetgreen!80!black] at (10.28, 3.05)
    {Synchronous Standby (AZ-2)};

  \node[stage, fill=quotientorange!15] (swrite) at (7.70, 2.12)
    {3. OS Cache\\ {\tiny \texttt{write\_lsn}}};
  \node[stage, fill=targetgreen!18] (sflush) at (10.28, 2.12)
    {4. Standby NVMe\\ {\tiny \texttt{flush\_lsn}}};
  \node[stage, fill=structurepurple!18] (sapply) at (12.86, 2.12)
    {5. Page Replay\\ {\tiny \texttt{replay\_lsn}}};

  \draw[->, very thick, sourceblue!85!black] (pdisk.east) -- node[above, font=\tiny]{TCP} (swrite.west);
  \draw[->, thick] (swrite) -- (sflush);
  \draw[->, thick] (sflush) -- (sapply);

  % Wait rungs below (all strictly within x = 0.25 .. 14.1)
  \draw[|<->|, thick, trapred] (0.25, 0.90) -- (1.40, 0.90);
  \node[rung, text=trapred] at (1.55, 0.90)
    {\texttt{off}: returns before local \texttt{fsync} (loses recent commits on OS crash)};

  \draw[|<->|, thick, sourceblue!85!black] (0.25, 0.42) -- (4.15, 0.42);
  \node[rung, text=sourceblue!85!black] at (4.30, 0.42)
    {\texttt{local}: local \texttt{fsync} only (RPO $> 0$ if AZ-1 disk dies)};

  \draw[|<->|, thick, quotientorange!90!black] (0.25, -0.06) -- (7.70, -0.06);
  \node[rung, text=quotientorange!85!black] at (7.85, -0.06)
    {\texttt{remote\_write}: standby OS cache ($\sim 1\text{ ms}$)};

  \draw[|<->|, thick, targetgreen!80!black] (0.25, -0.54) -- (10.28, -0.54);
  \node[rung, text=targetgreen!75!black] at (10.43, -0.54)
    {\texttt{on}: standby \texttt{fsync} ($\text{RPO}=0$)};

  \draw[|<->|, thick, structurepurple!90!black] (0.25, -1.02) -- (12.86, -1.02);
  \node[font=\scriptsize\bfseries, text=structurepurple!85!black, anchor=east] at (14.10, -1.42)
    {\texttt{remote\_apply}: replayed in standby \texttt{shared\_buffers} ($\text{RPO}=0$ + read-your-writes)};
\end{tikzpicture}
\caption{The five commit-wait watermarks of \texttt{synchronous\_commit} when \texttt{synchronous\_standby\_names} is active.}
\label{ch08:fig-replication-ladder}
\end{figure}

\begin{table}[H]
\centering
\small
\begin{tabular}{@{}llll@{}}
\toprule
\textbf{Mode} & \textbf{Wait Condition} & \textbf{Commit Latency} & \textbf{AZ-1 Loss RPO / Standby Reads} \\
\midrule
\texttt{off}          & WAL buffer copy       & $\sim 0.05\text{ ms}$ & Loses $\sim 200$--$600\text{ ms}$ on reboot \\
\texttt{local}        & Primary \texttt{fsync}& $\sim 0.3\text{ ms}$  & RPO $> 0$ on AZ loss; stale reads \\
\texttt{remote\_write}& Standby \texttt{write()}& $\sim 1.0\text{ ms}$& $\text{RPO}=0$ unless dual crash; brief lag \\
\texttt{on}           & Standby \texttt{fsync}& $\sim 1.8\text{ ms}$  & \textbf{$\text{RPO}=0$ on AZ loss}; brief lag \\
\texttt{remote\_apply}& Standby WAL replay    & $\sim 2.4\text{ ms}$  & \textbf{$\text{RPO}=0$ + causal standby reads} \\
\bottomrule
\end{tabular}
\caption{Durability (Recovery Point Objective, RPO) and intra-region cross-AZ latency across \texttt{synchronous\_commit} modes.}
\label{ch08:tab-sync-commit}
\end{table}

Two production lessons follow from Table~\ref{ch08:tab-sync-commit}. First, if
you leave \texttt{synchronous\_standby\_names = ''} (empty, the PostgreSQL
default), \texttt{synchronous\_commit = on} silently degrades to \texttt{local}:
PostgreSQL flushes only the primary's disk and ships WAL asynchronously. Second, in
a three-AZ cluster using \texttt{ANY 1 (az2\_standby, az3\_standby)}, setting
\texttt{synchronous\_commit = on} gives zero-data-loss (\texttt{RPO = 0}) survival
of an entire availability zone while tolerating one slow or rebooting standby
without halting writes; handlers that immediately redirect a customer to a
read-replica endpoint can elevate just that transaction via \texttt{SET LOCAL
synchronous\_commit = 'remote\_apply'}.

\subsection{Logical Replication for Zero-Downtime Upgrades and Event Streaming}
\label{ch08:logical-replication}

Physical replication requires the primary and standby to share the \emph{exact}
same major PostgreSQL version and CPU architecture because 8~KB page layouts and
catalog formats evolve across major releases. How, then, do you upgrade a live
financial ledger from PostgreSQL~16 to PostgreSQL~18 without taking a maintenance
outage, or stream committed \texttt{postings} rows into an analytics warehouse or
transactional outbox consumer?

Setting \texttt{wal\_level = logical} enables \textbf{logical decoding}. Instead of
shipping raw 8~KB block deltas, a logical \texttt{walsender} feeds WAL into the
\textbf{Reorder Buffer}, which reassembles interleaved physical WAL records per
\texttt{TransactionId} (spilling large transactions to disk if needed), discards
rolled-back transactions, and invokes an output plugin (\texttt{pgoutput}) the
moment a transaction commits to emit structured tuple-level \texttt{INSERT},
\texttt{UPDATE}, and \texttt{DELETE} messages:

\begin{lstlisting}[language=SQL]
-- On the primary (or PG 16 source cluster):
CREATE PUBLICATION ledger_pub
    FOR TABLE accounts, journal_entries, postings
    WITH (publish_via_partition_root = true);

-- On the PG 18 target cluster:
CREATE SUBSCRIPTION ledger_sub
    CONNECTION 'host=pg16-primary dbname=ledger user=repl_role'
    PUBLICATION ledger_pub;
\end{lstlisting}

Notice \texttt{publish\_via\_partition\_root = true}: logical subscribers receive
changes tagged under the parent table \texttt{postings}, allowing the target
cluster to use a completely different partitioning scheme or index set. Keep two
operational constraints in mind during logical upgrades: DDL statements
(\texttt{ALTER TABLE}) are not replicated automatically, and \texttt{SEQUENCE}
counters are not synchronized in real time---another reason time-ordered
\texttt{UUIDv7} primary keys simplify zero-downtime cutovers.

\section{Connection Pooling at Scale: Process Architecture and PgBouncer}
\label{ch08:pooling}

Our Rust \texttt{Axum} service runs on Tokio's work-stealing async runtime, where a
single pod handles ten thousand concurrent HTTP sockets in a few dozen megabytes of
RAM. Suppose you autoscale to $50$ Kubernetes pods, set each pod's
\texttt{sqlx::PgPool} \texttt{max\_connections} to $100$, and raise PostgreSQL's
\texttt{max\_connections = 5500} on a 32-core database server so all $5{,}000$
pooled sockets can stay open directly against PostgreSQL. Everything looks idle in
staging. Then Monday morning traffic arrives: database CPU spikes to $100\%$ inside
kernel spinlocks, free RAM vanishes, and P99 transfer latency climbs from $3\text{
ms}$ to four seconds.

\subsection{Why Five Thousand Backend Processes Melt PostgreSQL}
\label{ch08:process-model}

Unlike Tokio's cooperative async tasks, PostgreSQL uses a classic
\textbf{one-OS-process-per-connection} architecture: when a client opens a TCP
socket, the \texttt{postmaster} supervisor forks a dedicated \texttt{postgres}
backend process. Holding $5{,}000$ backend processes open degrades throughput
through three physical bottlenecks:

\begin{enumerate}
  \item \textbf{Per-backend private memory.} While \texttt{shared\_buffers} is
    shared across all processes, every \texttt{postgres} backend maintains its own
    private relation and system-catalog caches (\texttt{relcache} and
    \texttt{catcache}, which grow substantially on partitioned schemas), cached
    execution plans, and per-node \texttt{work\_mem} allocations (Chapter~7). At a
    realistic $10$--$16\text{ MB}$ RSS per active backend, $5{,}000$ processes
    consume $50$--$80\text{ GB}$ of RAM before caching a single heap page.
  \item \textbf{\texttt{ProcArray} MVCC snapshot scans.} Recall from Chapter~5 that
    every statement under \texttt{READ COMMITTED} calls \texttt{GetSnapshotData()},
    iterating across the shared-memory \texttt{ProcArray} of connected backends to
    record in-flight transaction IDs. Scanning $5{,}000$ backend slots hundreds of
    thousands of times per second causes heavy CPU cache-line bouncing around
    \texttt{ProcArrayLock} whenever commits update those same slots.
  \item \textbf{OS scheduler preemption inside locks.} When $600$ processes become
    runnable simultaneously on $32$ CPU cores, the Linux scheduler preempts
    backends every few milliseconds, flushing TLBs and L1/L2 caches, and
    stalling dozens of peers whenever a backend is context-switched
    off CPU \emph{while holding a lightweight buffer lock or spinlock}.
\end{enumerate}

A common starting point, which you should confirm by benchmarking your own
workload, is that peak transactional throughput on a server with
$N_{\text{cores}}$ CPU cores and NVMe storage occurs with roughly
$1.5 \times N_{\text{cores}}$ to $3 \times N_{\text{cores}}$ \emph{active}
backends: about \textbf{$48$ to $96$ backends on a 32-core primary}. Beyond that,
extra connections add contention without adding throughput.

\subsection{\texttt{PgBouncer} Transaction Multiplexing and the Prepared Statement Trap}
\label{ch08:pgbouncer}

How do $50$ Rust pods holding $2{,}500$ client connections share just $48$ physical
\texttt{postgres} processes? You place \textbf{\texttt{PgBouncer}}---a lightweight,
event-driven C proxy that speaks the PostgreSQL v3 wire protocol at roughly
$2\text{ KB}$ of memory per client socket---in front of the database
(Figure~\ref{ch08:fig-pgbouncer}). \texttt{PgBouncer} offers three pooling modes:

\begin{itemize}
  \item \textbf{\texttt{session} mode} pins a physical \texttt{postgres} backend to
    a client socket from connect until disconnect. Because \texttt{sqlx::PgPool}
    holds long-lived TCP connections open, \texttt{session} mode achieves \emph{zero}
    multiplexing ($2{,}500$ client sockets still demand $2{,}500$ backends).
  \item \textbf{\texttt{transaction} mode} (the production standard) leases a
    physical backend to a client \emph{only from \texttt{BEGIN} (or a standalone
    query) until \texttt{COMMIT} or \texttt{ROLLBACK}}. The instant \texttt{COMMIT}
    completes, \texttt{PgBouncer} returns that backend to the warm server pool for
    the next waiting request. While an \texttt{Axum} handler parses JSON, checks
    cryptographic signatures, or streams HTTP responses, it occupies \emph{zero}
    PostgreSQL processes.
  \item \textbf{\texttt{statement} mode} reclaims the backend after every single
    SQL statement and immediately aborts any multi-statement \texttt{BEGIN ...
    COMMIT} block---unusable for double-entry transactions.
\end{itemize}

\begin{figure}[H]
\centering
\begin{tikzpicture}[
  font=\small,
  >=Latex,
  pod/.style={draw=sourceblue!85!black, fill=sourceblue!10, rounded corners=3pt,
              minimum width=3.2cm, minimum height=0.85cm, align=center, font=\scriptsize},
  bouncer/.style={draw=quotientorange!90!black, fill=quotientorange!14, thick, rounded corners=4pt,
                  minimum width=3.9cm, minimum height=2.9cm, align=center, font=\scriptsize},
  backend/.style={draw=targetgreen!80!black, fill=targetgreen!14, rounded corners=2pt,
                  minimum width=2.8cm, minimum height=0.62cm, align=center, font=\scriptsize}
]
  % Left: 50 Rust Axum pods (x: 0 .. 3.4)
  \node[font=\small\bfseries, text=sourceblue!85!black] at (1.7, 3.35)
    {50 Rust \texttt{Axum} Pods};
  \node[pod] (p1) at (1.7, 2.50) {Pod 1: Tokio tasks\\ {\tiny \texttt{PgPool} (\texttt{max=50})}};
  \node[pod] (p2) at (1.7, 1.40) {Pod 2: Tokio tasks\\ {\tiny \texttt{PgPool} (\texttt{max=50})}};
  \node[font=\bfseries] at (1.7, 0.72) {$\vdots$};
  \node[pod] (pn) at (1.7, 0.05) {Pod 50: Tokio tasks\\ {\tiny \texttt{PgPool} (\texttt{max=50})}};

  % Center: PgBouncer (x: 5.1 .. 9.0)
  \node[bouncer] (pgb) at (7.05, 1.30) {
    \textbf{\texttt{PgBouncer} ($\ge 1.21$)}\\[3pt]
    \texttt{pool\_mode = transaction}\\
    \texttt{max\_client\_conn = 5000}\\
    \texttt{default\_pool\_size = 48}\\
    \texttt{max\_prepared\_statements}\\
    \texttt{= 200}\\[2pt]
    {\tiny Maps client \texttt{sqlx\_s\_1}}\\
    {\tiny $\leftrightarrow$ backend \texttt{PGBouncer\_7}}
  };

  % Right: PostgreSQL Primary (x: 10.6 .. 14.3)
  \draw[rounded corners=4pt, draw=targetgreen!80!black, fill=targetgreen!4]
    (10.55, -0.55) rectangle (14.35, 3.15);
  \node[font=\small\bfseries, text=targetgreen!80!black] at (12.45, 2.85)
    {PostgreSQL (32 Cores)};

  \node[backend] (b1) at (12.45, 2.15) {\texttt{postgres} \#1 (active)};
  \node[backend] (b2) at (12.45, 1.35) {\texttt{postgres} \#2 (active)};
  \node[backend] (b3) at (12.45, 0.55) {\texttt{postgres} \#48 (active)};
  \node[draw=structurepurple, fill=structurepurple!14, rounded corners=2pt,
        minimum width=3.3cm, minimum height=0.58cm, font=\scriptsize]
    at (12.45, -0.15) {Compact \texttt{ProcArray} + \texttt{shared\_buffers}};

  % Arrows
  \draw[->, thick, sourceblue!80!black] (p1.east) -- node[above, sloped, font=\tiny]{2,500 TCP} ([yshift=22pt]pgb.west);
  \draw[->, thick, sourceblue!80!black] (p2.east) -- ([yshift=2pt]pgb.west);
  \draw[->, thick, sourceblue!80!black] (pn.east) -- ([yshift=-24pt]pgb.west);

  \draw[->, very thick, targetgreen!80!black] ([yshift=24pt]pgb.east) -- node[above, font=\tiny]{48 warm} (b1.west);
  \draw[->, very thick, targetgreen!80!black] ([yshift=2pt]pgb.east) -- (b2.west);
  \draw[->, very thick, targetgreen!80!black] ([yshift=-20pt]pgb.east) -- (b3.west);
\end{tikzpicture}
\caption{Multiplexing $2{,}500$ \texttt{SQLx} client connections across $48$ warm \texttt{postgres} backends via \texttt{PgBouncer} in \texttt{transaction} mode, keeping \texttt{ProcArray} compact and CPU caches hot.}
\label{ch08:fig-pgbouncer}
\end{figure}

Switching \texttt{pool\_mode = transaction} without understanding the PostgreSQL
wire protocol triggers a classic production outage within seconds:
\texttt{ERROR: prepared statement "sqlx\_s\_1" does not exist} (or
\texttt{already exists}). Here is the exact packet-level reason:

\begin{enumerate}
  \item By default, \texttt{SQLx} uses PostgreSQL's extended query protocol and
    caches up to $100$ named prepared statements per TCP socket. In Transaction~A,
    Client~1 is assigned \texttt{postgres} Backend~\#4 and sends a \texttt{Parse}
    message creating statement \texttt{sqlx\_s\_1} inside Backend~\#4's private
    process memory, then commits.
  \item One millisecond later, Client~1 starts Transaction~B and \texttt{PgBouncer}
    routes it to Backend~\#19. Assuming its TCP socket still talks to the same
    process, \texttt{SQLx} skips \texttt{Parse} and sends a \texttt{Bind} packet
    naming \texttt{sqlx\_s\_1}. Backend~\#19 has never heard of \texttt{sqlx\_s\_1}
    and aborts the transaction.
\end{enumerate}

The same backend-hopping mechanism breaks any state stored at \emph{session} scope
rather than \emph{transaction} scope:
\begin{itemize}
  \item Never run bare \texttt{SET synchronous\_commit = 'remote\_apply'} or
    \texttt{SET statement\_timeout = '2s'} in transaction mode---the setting leaks
    onto the next unrelated request that lands on that backend. Always write
    \textbf{\texttt{SET LOCAL ...}}, which PostgreSQL automatically reverts at
    \texttt{COMMIT} or \texttt{ROLLBACK}.
  \item Never call session-scoped \texttt{pg\_advisory\_lock}, which leaves the lock
    held on a pooled backend after your handler returns; always use the
    transaction-scoped \textbf{\texttt{pg\_advisory\_xact\_lock}} function instead.
  \item \texttt{LISTEN / NOTIFY} pub-sub channels require a dedicated direct or
    \texttt{session}-mode connection outside the transaction pool.
\end{itemize}

For prepared statements, you have two clean fixes. Best of all, \textbf{modern
\texttt{PgBouncer} ($\ge 1.21$)} supports protocol-level prepared-statement
multiplexing when you set \texttt{max\_prepared\_statements = 200} in
\texttt{pgbouncer.ini}: \texttt{PgBouncer} tracks each client's \texttt{Parse}
text, rewrites client names (\texttt{sqlx\_s\_1}) into deduplicated server-side
names (\texttt{PGBOUNCER\_7}), and transparently prepares any missing statement on
whichever physical backend leases the transaction---preserving Chapter~7's plan-cache
speedups. When deploying behind older proxies or cloud poolers without protocol
tracking, disable client-side named statement caching on \texttt{PgConnectOptions}
so \texttt{SQLx} sends unnamed single-use statements:

\begin{lstlisting}[language=Rust]
use sqlx::postgres::{PgConnectOptions, PgPoolOptions};
use std::str::FromStr;
use std::time::Duration;

pub async fn build_ledger_pool(
    db_url: &str,
    pgbouncer_tracks_stmts: bool,
) -> Result<sqlx::PgPool, sqlx::Error> {
    let mut opts = PgConnectOptions::from_str(db_url)?;
    if !pgbouncer_tracks_stmts {
        // Fallback for stateless transaction poolers: unnamed statements.
        opts = opts.statement_cache_capacity(0);
    }

    PgPoolOptions::new()
        .max_connections(50)
        .acquire_timeout(Duration::from_secs(2))
        .idle_timeout(Duration::from_secs(300))
        .connect_with(opts)
        .await
}
\end{lstlisting}

Here is the whole path of one \$50.00 transfer through the machinery of this
book. An \texttt{Axum} task borrows a pooled connection through \texttt{PgBouncer}
and lands on one of 48 backends, so snapshot scans of \texttt{ProcArray} stay
short (Chapter~5). It inserts rows with \texttt{UUIDv7} keys, which append to the
rightmost B-tree leaf already in \texttt{shared\_buffers} (Chapters~2 and~3). It
claims its idempotency key through a \texttt{UNIQUE} index, locks both accounts in
ascending order with \texttt{FOR NO KEY UPDATE} (Chapter~6), and inserts two
\texttt{BIGINT} postings into the current month's partition. At \texttt{COMMIT},
the deferred trigger checks that the legs sum to zero, and with
\texttt{synchronous\_commit = on} the commit record is flushed to WAL on two
machines before the client hears success (Chapter~4).

\section{Exercises}
\label{ch08:exercises}

\begin{exercise}[IEEE 754 Drift vs.\ Exact Ledger Representations]
Execute the following query in PostgreSQL:
\begin{lstlisting}[language=SQL]
SELECT SUM(val::float8) AS float_sum, SUM(val::numeric) AS num_sum
FROM generate_series(1, 1000000) AS g(i),
     LATERAL (VALUES (0.10), (0.20), (-0.30)) AS v(val);
\end{lstlisting}
Explain why \texttt{float\_sum} is nonzero while \texttt{num\_sum} is exact. Next,
design the columns for an FX conversion table that converts \texttt{JPY} (scale $0$),
\texttt{USD} (scale $2$), and \texttt{BTC} (scale $8$) using \texttt{BIGINT} minor
units for posted legs and \texttt{NUMERIC(24, 12)} for the exchange rate,
specifying where rounding happens and which rule it uses so the converted
journal entry still obeys conservation within each currency. (Note that
PostgreSQL's \texttt{round()} on \texttt{NUMERIC} rounds halves away from zero,
not to even; banker's rounding needs an explicit expression.)
\end{exercise}

\begin{exercise}[The Autumn Daylight-Saving Collision]
In a \texttt{psql} session with \texttt{TimeZone} set to \texttt{'America/New\_York'},
compare storing \texttt{'2026-11-01 01:30:00-04'} and \texttt{'2026-11-01
01:30:00-05'} in a \texttt{TIMESTAMP} column versus a \texttt{TIMESTAMPTZ} column.
Compute the exact difference between the two 64-bit integer values stored on the
8~KB heap page for \texttt{TIMESTAMPTZ} (in microseconds since
\texttt{2000-01-01 00:00:00 UTC}), and show how \texttt{TIMESTAMP} corrupts an
hourly TWAP (time-weighted average price) query across that hour.
\end{exercise}

\begin{exercise}[Quantifying \texttt{UUIDv4} WAL Amplification vs.\ \texttt{UUIDv7}]
Suppose a primary-key B-tree on \texttt{postings} has $8{,}000{,}000$ leaf pages
($\approx 64\text{ GB}$). Immediately after a \texttt{CHECKPOINT} completes (with
\texttt{full\_page\_writes = on}), a burst inserts $40{,}000$ new rows within two
seconds. Estimate how many distinct 8~KB leaf pages are dirtied---and how many
megabytes of 8~KB Full Page Images are written to WAL---under random \texttt{UUIDv4}
versus time-ordered \texttt{UUIDv7} (assuming $\sim 120$ index tuples fit per new
rightmost leaf page).
\end{exercise}

\begin{exercise}[Non-Overlapping Credit Limits with \texttt{EXCLUDE USING gist}]
Write the DDL for a \texttt{credit\_limits} table where each \texttt{account\_id}
has an overdraft limit valid over an \texttt{effective\_during} column of type
\texttt{TSTZRANGE}, ensuring that \emph{active} schedules (\texttt{status =
'active'}) for the same account can never overlap in time while allowing historical
\texttt{status = 'superseded'} rows to remain in the table for audit purposes.
(Hint: inspect PostgreSQL's partial exclusion constraint syntax \texttt{EXCLUDE
USING gist (...) WHERE (...)}.) Test what happens at the boundary when two
intervals use closed bounds \texttt{'[]'} instead of half-open \texttt{'[)'}.
\end{exercise}

\begin{exercise}[Auditing and Hardening the Double-Entry Deferred Trigger]
Suppose a transaction inserts four postings for the same \texttt{entry\_id} (sender
debit, recipient credit, platform fee, partner rebate). How many times does a
\texttt{FOR EACH ROW} deferred constraint trigger fire at \texttt{COMMIT}? Show how
to deduplicate redundant checks within the same transaction (for instance, skipping
if a per-transaction temporary table or session variable already verified
\texttt{v\_entry\_id} at the final command counter), and write the companion
deferred constraint trigger on \texttt{journal\_entries} that rejects a header
committed with zero rows in \texttt{postings}.
\end{exercise}

\begin{exercise}[Deadlock Cycles and \texttt{FOR NO KEY UPDATE} Internals]
Construct a two-session \texttt{psql} interleaving ($T_1$ transferring $A \to B$
and $T_2$ transferring $B \to A$) that deterministically triggers
\texttt{ERROR: deadlock detected} when canonical UUID sorting is omitted from
Listing~\ref{ch08:lst-rust-transfer}. Then, using Chapter~6's tuple header lock
modes, demonstrate in two \texttt{psql} sessions why locking \texttt{accounts} with
\texttt{FOR UPDATE} blocks a concurrent \texttt{INSERT INTO postings} foreign-key
check, whereas \texttt{FOR NO KEY UPDATE} lets the foreign-key check proceed until
the cached-balance \texttt{UPDATE} step.
\end{exercise}

\begin{exercise}[Partitioning and Retention on a Multi-Billion-Row Table]
Write a SQL migration or PL/pgSQL helper procedure \texttt{ensure\_monthly\_partition}
taking \texttt{target\_month DATE} that safely creates next month's
\texttt{postings\_YYYY\_MM} range partition ahead of time. Explain what happens to
live \texttt{INSERT} traffic at \texttt{2027-01-01 00:00:00Z} if the cron job fails
to create January's partition and no \texttt{DEFAULT} partition exists---and
conversely, why adding a bloated \texttt{DEFAULT} partition can turn a future
\texttt{ATTACH PARTITION} into a long table-scan lock.
\end{exercise}

\begin{exercise}[Sizing \texttt{synchronous\_commit} Across Three Availability Zones]
A ledger cluster spans \texttt{az1} (primary), \texttt{az2} (standby), and
\texttt{az3} (standby) with local NVMe \texttt{fsync} latency $0.25\text{ ms}$,
round-trip network latency $1.1\text{ ms}$ (\texttt{az1}--\texttt{az2}) and
$1.6\text{ ms}$ (\texttt{az1}--\texttt{az3}), and standby WAL replay overhead
$0.35\text{ ms}$. With \texttt{synchronous\_standby\_names = 'ANY 1 (az2, az3)'},
calculate the expected commit latency under normal operation and during an
\texttt{az2} outage for \texttt{local}, \texttt{remote\_write}, \texttt{on}, and
\texttt{remote\_apply}. Which mode should \texttt{POST /v1/transfers} use if
\texttt{GET /v1/accounts/:id/balance} is load-balanced across \texttt{az2} and
\texttt{az3}?
\end{exercise}

\begin{exercise}[Fixing a \texttt{PgBouncer} Transaction-Mode Rollout]
An engineering team flips \texttt{PgBouncer} from \texttt{pool\_mode = session} to
\texttt{pool\_mode = transaction} in front of $40$ Rust \texttt{Axum} pods. Within
two minutes, monitoring lights up with three errors: (a) \texttt{prepared statement
``sqlx\_s\_4'' does not exist}, (b) fast transfers mysteriously timing out at
$150\text{ ms}$ after an admin endpoint ran \texttt{SET statement\_timeout =
'150ms'}, and (c) a nightly settlement job holding \texttt{pg\_advisory\_lock(99)}
never releasing the lock after commit. Trace the exact backend-multiplexing cause
of all three bugs and write the complete configuration and SQL/Rust fixes.
\end{exercise}
